%%%%%%%%%%%%%%%%%%%%%%%%%%%%%%%%%%%%%%%%
%
% $ Autor: Gruppe 07 $
% $ Projekt: Magic Faucet Fountain $
% $ Pfad: LaTeXTemplate/Nano33BLESense/Handbuch Magic Faucet Fountain/Chapters/de/Skizze.tex $
% $ Kurzbeschreibung: Dieses Dokument zeigt die Skizze der Applikation. $
%
% !TeX spellcheck = de_DE/GB
% !TeX program = pdflatex
% !BIB program = biber/bibtex
% !TeX encoding = utf8
%
%%%%%%%%%%%%%%%%%%%%%%%%%%%%%%%%%%%%%%%%

\chapter{Skizze}

\section{Übersichtsskizze Blumentopf und Netzteilkasten}
\begin{figure}[!h]
	\centering
	\resizebox{1\textwidth}{!}{%
		\begin{circuitikz}
			\tikzstyle{every node}=[font=\Huge]
			\draw [line width=1pt, short] (0,5.75) -- (11.25,5.75);
			\draw [line width=1pt, short] (11.25,5.75) -- (15,20.75);
			\draw [line width=1pt, short] (0,5.75) -- (-3.75,20.75);
			\draw [line width=1pt, short] (-3.75,20.75) -- (-5,20.75);
			\draw [line width=1pt, short] (-5,20.75) -- (-5,21.5);
			\draw [line width=1pt, short] (-5,21.5) -- (-3,21.5);
			\draw [line width=1pt, short] (-3,21.5) -- (0.75,6.5);
			\draw [line width=1pt, short] (0.75,6.5) -- (10.5,6.5);
			\draw [line width=1pt, short] (15,20.75) -- (16.25,20.75);
			\draw [line width=1pt, short] (16.25,20.75) -- (16.25,21.5);
			\draw [line width=1pt, short] (16.25,21.5) -- (14.25,21.5);
			\draw [line width=1pt, short] (14.25,21.5) -- (10.5,6.5);
			\draw [line width=1pt, short] (-5,21.5) -- (-5,22);
			\draw [line width=1pt, short] (16.25,22) -- (16.25,21.5);
			\draw [line width=1pt, short] (14.25,21.5) -- (11.25,15.75);
			\draw [line width=1pt, short] (-3,21.5) -- (0,15.75);
			\draw [line width=1pt, short] (-2.5,22) -- (0.5,16.25);
			\draw [line width=1pt, short] (10.75,16.25) -- (13.75,22);
			\draw [ line width=1pt ] (7.75,15.75) rectangle (10,15);
			\draw [ line width=1pt ] (8,14.75) rectangle (8.75,14.5);
			\draw [ line width=1pt ] (9,14.75) rectangle (9.75,14.5);
			\draw [line width=1pt, short] (5.75,7.5) -- (5.75,30);
			\draw [line width=1pt, short] (6,7.5) -- (6,30);
			\draw [ line width=1pt ] (5.75,8) rectangle (6,7.5);
			\draw [ line width=1pt ] (4.75,7.5) rectangle (6,6.75);
			\draw [ line width=1pt ] (4.5,7) rectangle (4.75,6.75);
			\draw [ line width=1pt ] (8,15) rectangle (9.75,14.75);
			\draw [line width=1pt, short] (7.75,15) -- (7.75,12.5);
			\draw [line width=1pt, short] (10,15) -- (10,12.5);
		%	\draw [line width=1pt, short] (7.75,8.25) -- (8,8.25);
		%	\draw [line width=1pt, short] (10,8.25) -- (9.75,8.25);
		%	\draw [line width=1pt, short] (9.75,8.25) -- (9.75,7);
		%	\draw [line width=1pt, short] (8,8.25) -- (8,7);
		%	\draw [line width=1pt, short] (8,7) -- (9.75,7);
			\draw [line width=1pt, short] (7.5,15.75) -- (7.5,12.5);
			\draw [line width=1pt, short] (10.25,15.75) -- (10.25,12.5);
			\draw [line width=1pt, short] (7.5,12.5) -- (10.25,12.5);
		%	\draw [line width=1pt, short] (10.25,7.5) -- (10.25,6.75);
		%	\draw [line width=1pt, short] (10.25,6.75) -- (7.5,6.75);
		%	\draw [line width=1pt, short] (7.75,7.5) -- (7.5,7.5);
		%	\draw [line width=1pt, short] (7.5,7.5) -- (7.5,6.75);
		%	\draw [ line width=1pt ] (8,9) rectangle (9.75,8.5);
			\draw [line width=1pt, short] (10,14.25) -- (10.25,14.25);
			\draw [line width=1pt, short] (10,14) -- (10.25,14);
			\draw [line width=1pt, short] (12.5,14) -- (13.25,14);
			\draw [line width=1pt, short] (12.25,13.25) -- (13,13.25);
			\draw [line width=1pt, short] (0,15.75) -- (5.5,15.75);
			\draw [line width=1pt, short] (6.25,15.75) -- (11.25,15.75);
			\draw [line width=1pt, short] (10.75,16.25) -- (6.25,16.25);
			\draw [line width=1pt, short] (0.5,16.25) -- (5.5,16.25);
			\draw [line width=1pt, short] (-5,22) -- (5.5,22);
			\draw [line width=1pt, short] (6.25,22) -- (16.25,22);
			\draw [ line width=1pt ] (21.25,8.25) rectangle (26.75,5.75);
			\draw [line width=1pt, short] (21.25,7) -- (26.75,7);
			\draw [line width=1pt, short] (21.25,7) -- (20.25,7);
			\draw [line width=1pt, short] (25,8.25) -- (25,8.75);
			\draw [line width=1pt, short] (25,8.75) -- (25.5,8.5);
			\draw [line width=1pt, short] (26,8.75) -- (26,8.25);
			\draw [line width=1pt, short] (25.5,8.5) -- (26,8.75);
			\draw [line width=1pt, short] (13.25,13.5) -- (14,13.5);
			\draw [ line width=1pt ] (5.25,30.5) rectangle (6.5,30);
			\draw [line width=1pt, short] (5.5,30.5) .. controls (5.5,31.75) and (5.75,32.25) .. (7,32.5);
			\draw [line width=1pt, short] (6.25,30.5) .. controls (6.25,31.5) and (6.5,31.5) .. (7.5,31.75);
			\draw [line width=1pt, short] (7,32.5) .. controls (7.25,32.5) and (7.5,32.5) .. (7.5,32.75);
			\draw [line width=1pt, short] (7.5,32.75) -- (7.5,33.25);
			\draw [line width=1pt, short] (7.5,33.25) -- (8.75,33.25);
			\draw [line width=1pt, short] (8.75,33.25) -- (8.75,32.75);
			\draw [line width=1pt, short] (7.5,31.75) -- (9.5,31.75);
			\draw [line width=1pt, short] (8.75,32.75) .. controls (8.75,32.75) and (8.75,32.5) .. (9,32.5);
			\draw [line width=1pt, short] (9,32.5) -- (9.5,32.5);
			\draw [ line width=1pt ] (9.5,32.75) rectangle (9.75,31.5);
			\draw [ line width=1pt ] (8,33.25) rectangle (8.25,33.75);
			\draw [ line width=1pt ] (7.5,34) rectangle (8.75,33.75);
			\draw [ line width=1pt ] (6,29.75) rectangle (6.25,29.5);
			\draw [ line width=1pt ] (5.5,29.75) rectangle (5.75,29.5);
			\draw [line width=1pt, ->, >=Stealth] (17.5,27) -- (12.5,22.25);
			\draw [line width=1pt, ->, >=Stealth] (29,13) -- (25.5,8.75);
			\node [font=\Huge] at (21.75,27.75) {Abnehmbarer Deckel};
			\node [font=\Huge] at (32,13.75) {Hauptschalter};
			\draw [ line width=1pt ] (5.5,30) rectangle (6.25,7.5);
		\end{circuitikz}
	}%
\label{fig:Skizze}
\caption{Magic Faucet Fountain}
\end{figure}

\newpage

%\section{Aufbau des Netzteilkastens}
%\begin{figure}[!ht]
%	\centering
%	\resizebox{0.7\textwidth}{!}{%
%		\begin{circuitikz}
%			\tikzstyle{every node}=[font=\small]
%			\draw [ line width=1pt](2.5,12) to[short] (11.25,12);
%			\draw [ line width=1pt](11.25,12) to[short] (11.25,5.75);
%			\draw [ line width=1pt](11.25,5.75) to[short] (2.5,5.75);
%			\draw [ line width=1pt](2.5,5.75) to[short] (2.5,12);
%			\draw [ line width=1pt](5,10.75) to[short] (2.5,10.75);
%			\draw [ line width=1pt](5,10.75) to[short] (5,7);
%			\draw [ line width=1pt](5,7) to[short] (2.5,7);
%			\draw [line width=1pt, short] (6.25,10.75) -- (10,10.75);
%			\draw [line width=1pt, short] (10,10.75) -- (10,9.5);
%			\draw [line width=1pt, short] (10,9.5) -- (6.25,9.5);
%			\draw [line width=1pt, short] (6.25,9.5) -- (6.25,10.75);
%			\draw [ line width=1pt](4.25,10) to[short, -o] (5.25,10) ;
%			\draw [ line width=1pt](4,7.75) to[short, -o] (5.5,7.75) ;
%			\draw [ line width=1pt](2.75,9.5) to[short, -o] (1.25,9.5) ;
%			\draw [ line width=1pt](2.75,8.75) to[short, -o] (1.25,8.75) ;
%			\draw [ line width=1pt](2.75,8) to[short, -o] (1.25,8) ;
%			\node [font=\normalsize, color={rgb,255:red,255; green,0; blue,0}] at (0.25,9.5) {230 V};
%			\node [font=\normalsize, color={rgb,255:red,0; green,128; blue,255}] at (0.5,8.75) {N};
%			\node [font=\normalsize, color={rgb,255:red,0; green,255; blue,64}] at (0.5,8) {PE};
%			\draw [ line width=1pt](4,8.75) to[short, -o] (5.5,8.75) ;
%			\node [font=\normalsize] at (4.25,7.5) {24 V};
%			\node [font=\normalsize] at (4.25,8.5) {GND};
%			\node [font=\normalsize] at (4.25,9.75) {5 V};
%			\node [font=\normalsize] at (4,11) {\textbf{Schaltnetzteil}};
%			\node [font=\normalsize] at (7.75,11) {\textbf{Arduino}};
%			\draw [ line width=1pt](6.5,10) to[short, -o] (6,10) ;
%			\draw [ line width=1pt](7,9.75) to[short, -o] (7,9.25) ;
%			\draw [ color={rgb,255:red,255; green,0; blue,0}, , line width=1pt](5.25,10) to[short] (6,10);
%			\draw [ color={rgb,255:red,128; green,64; blue,64}, , line width=1pt](5.5,8.75) to[short] (7,8.75);
%			\draw [ color={rgb,255:red,128; green,64; blue,64}, , line width=1pt](7,9.25) to[short] (7,8.75);
%			\draw [line width=1pt](8,7.75) to[Tnpn, transistors/scale=1.19] (10,7.75);
%			\draw [ line width=1pt](9,9.75) to[short, -o] (9,9.25) ;
%			\draw [ color={rgb,255:red,255; green,0; blue,0}, , line width=1pt](5.5,7.75) to[short] (8,7.75);
%			\draw [ color={rgb,255:red,0; green,255; blue,0}, , line width=1pt](9,9.25) to[short] (9,8.75);
%			\draw [ color={rgb,255:red,255; green,0; blue,0}, , line width=1pt](10,7.75) to[short] (12,7.75);
%			\node [font=\normalsize] at (12.75,7.75) {\textbf{Pumpe}};
%			\node [font=\normalsize] at (9,7.25) {\textbf{MOSFET}};
%			\node [font=\small] at (7.5,9.75) {GND};
%			\node [font=\small] at (9.5,9.75) {DO1};
%			\node [font=\small] at (6.75,10.25) {3,3 V};
%		\end{circuitikz}
%	}%
%	\caption{Der Netzteilkasten}
%	\label{fig:netzteilkasten}
%\end{figure}

%\section{Aufbau des Schwimmerrohres}
%\begin{figure}[!ht]
%	\centering
%	\resizebox{0.6\textwidth}{!}{%
%		\begin{circuitikz}
%			\tikzstyle{every node}=[font=\LARGE]
%			\draw [short] (5,3.25) -- (5,13.25);
%			\draw [short] (8.75,3.25) -- (8.75,13.25);
%			\draw [short] (5,13.25) -- (8.75,13.25);
%			\draw [short] (8.25,3.25) -- (8.25,12.25);
%			\draw [short] (5.5,3.25) -- (5.5,12);
%			\draw [short] (5.5,12) -- (5.5,12.25);
%			\draw [short] (5.5,12.25) -- (8.25,12.25);
%			\draw [ color={rgb,255:red,0; green,128; blue,255}, line width=0.6pt, short] (5.75,12.25) -- (5.75,11.5);
%			\draw [ color={rgb,255:red,0; green,128; blue,255}, line width=0.6pt, short] (5.75,11.5) -- (6.5,11.5);
%			\draw [ color={rgb,255:red,0; green,128; blue,255}, line width=0.6pt, short] (6.5,11.5) -- (6.5,12.25);
%			\draw [ color={rgb,255:red,0; green,128; blue,255}, line width=0.6pt, short] (8,12.25) -- (8,11.5);
%			\draw [ color={rgb,255:red,0; green,128; blue,255}, line width=0.6pt, short] (8,11.5) -- (7.25,11.5);
%			\draw [ color={rgb,255:red,0; green,128; blue,255}, line width=0.6pt, short] (7.25,11.5) -- (7.25,12.25);
%			\draw [ color={rgb,255:red,0; green,128; blue,255}, line width=0.6pt, short] (5.75,12.25) -- (6.5,12.25);
%			\draw [ color={rgb,255:red,0; green,128; blue,255}, line width=0.6pt, short] (7.25,12.25) -- (8,12.25);
%			\draw [ color={rgb,255:red,0; green,128; blue,255}, line width=0.6pt, short] (7.25,12) -- (6.5,12);
%			\begin{scope}[rotate around={-87.5:(6.75,11.25)}]
%				\draw[domain=6.75:12.25,samples=100,smooth color={rgb,255:red,255; green,0; blue,0}, , line width=0.6pt] plot (\x,{1*sin(6.22*\x r -6.75 r ) +11.25});
%			\end{scope}
%			\draw [ color={rgb,255:red,128; green,64; blue,0} , fill={rgb,255:red,128; green,64; blue,64}, line width=0.6pt ] (7,4.75) ellipse (1cm and 0.5cm);
%			\draw[domain=5.5:8.25,samples=100,smooth color={rgb,255:red,0; green,128; blue,255}, , line width=0.6pt] plot (\x,{0.1*sin(10*\x r -0.5 r ) +4});
%			\draw [line width=0.6pt, short] (5,3.25) -- (5.5,3.25);
%			\draw [line width=0.6pt, short] (8.25,3.25) -- (8.75,3.25);
%			\draw[domain=2.25:5,samples=100,smooth color={rgb,255:red,0; green,128; blue,255}, , line width=0.6pt] plot (\x,{0.1*sin(10*\x r +2.75 r ) +4});
%			\draw[domain=8.75:11.5,samples=100,smooth color={rgb,255:red,0; green,128; blue,255}, , line width=0.6pt] plot (\x,{0.1*sin(10*\x r -3.75 r ) +4});
%			\node [font=\LARGE] at (6,11.75) {R};
%			\node [font=\LARGE] at (7.75,11.75) {T};
%		\end{circuitikz}
%	}%
%	\caption{Aufbau und Funktionsweise vom Schwimmer-Rohr}
%	\label{fig:schwimmer}
%\end{figure}

